\section{Discussion}

\subsection{Implications for In-Situ Resource Utilization (ISRU) and Deep Space Agriculture}

Understanding the fine-scale physical, mineralogical, and environmental characteristics of extraterrestrial regolith is critical for Bioregenerative Life Support Systems (BLSS) and deep space agriculture~\cite{barker2023deepspaceag}. 

\begin{itemize}
    \item \textbf{Mineral Nutrition & Substrate Suitability}: The identification of high-iron, intermediate-titanium olivine-pyroxene assemblages in Mare Imbrium (Chang'e-3) versus deep-mantle low-Ca pyroxenes in the South Pole-Aitken Basin (Chang'e-4) highlights significant compositional heterogeneity. For lunar plant growth and regolith conditioning, high-titanium and fine glassy agglutinates can affect cation exchange capacity, microbial colonization, and root-zone water retention.
    \item \textbf{Water & Volatiles Distribution}: The in-situ verification of indigenous water and hydroxyl ($OH$) in Chang'e-5 young basalts ($2.85~\mu\text{m}$ band, up to 120~ppm in soil and 900~ppm in apatite) confirms that lunar volcanism preserved endogenous water up to 2.0 billion years ago. On Mars, the MarSCoDe detection of hydrated sulfates in Amazonian duricrust indicates accessible water-bearing minerals in Utopia Planitia for extraction and ISRU processing.
    \item \textbf{Subsurface Stability & Habitat Placement}: The GPR radargram stratigraphy from Chang'e-4 and Zhurong provides direct structural guidelines for subsurface drilling and shelter construction. The upper $12\text{ m}$ fine ejecta sheet on the Moon provides easily excavatable regolith with excellent thermal insulation and low dielectric loss, ideal for surface habitat berming.
\end{itemize}

\subsection{Space Environmental Stressors: Radiation and Hypomagnetic Fields}

Surface biology and future astronaut crews face co-occurring physical stressors. The continuous dose equivalent rate of $1.369\text{ mSv/day}$ measured by Chang'e-4 LND on the lunar farside is roughly 200 times higher than Earth's surface and 2.6 times higher than the International Space Station (ISS) interior, emphasizing the necessity of $>2\text{ m}$ regolith overburden for long-term habitats.

Furthermore, both the lunar surface and Martian northern lowlands (Utopia Planitia, RoMAG $\sim 100\text{ nT}$) exhibit extreme hypomagnetic field (HMF) conditions ($< 1~\mu\text{T}$, compared to Earth's $\sim 50~\mu\text{T}$). Prolonged HMF exposure is known to disrupt radical pair spin-state recombination kinetics in cryptochrome flavoproteins, impairing circadian rhythms, reactive oxygen species (ROS) regulation, and root gravitropism in plants. The integration of magnetic and radiation timeseries in our visualization suite facilitates holistic biological risk assessments.

\subsection{FAIR Principles and Open Planetary Data Science}

Planetary exploration data often suffer from accessibility barriers due to complex PDS binary formats and localized archiving portals. By transforming raw CLPDS datasets into standardized JSON schemas, modular JavaScript visualization engines, and open-source WebGL environments, this project provides a reproducible template for international planetary data democratization.
